\chapter{Agent \texorpdfstring{$\leftrightarrow$}{<->} Lean 4 Pipeline}
\label{ch:pipeline}

\emph{Agentic Hallucinations --- Independent Study.} Design spec for how the agent and Lean 4 actually communicate, and what ``environment switching'' means across this project's full stack. Written to be concrete enough to hand to an implementer (human or Claude Code) as a blueprint --- not just a description.

Sources: \href{https://github.com/leanprover-community/repl}{leanprover-community/repl} (protocol fetched and quoted directly below), \href{https://arxiv.org/abs/2410.16429}{Pantograph (arXiv:2410.16429)}, \href{https://github.com/oOo0oOo/lean-lsp-mcp}{lean-lsp-mcp}, \href{https://github.com/augustepoiroux/leaninteract}{LeanInteract}.

\section{The loop, restated precisely}

From the original project brief: the agent is given a theorem, proposes one tactic, Lean verifies it, and either returns a new goal state or an error --- repeat until the goal list is empty (proved) or some stopping condition is hit (budget exhausted, agent gives up, non-termination).

Stated as a state machine:

\begin{verbatim}
state = (theorem, goal_stack, history)
loop:
  agent.observe(goal_stack, history) → tactic_string
  result = lean.apply(tactic_string, goal_stack)
  if result.ok and result.goals == []:
      state = PROVED; stop
  elif result.ok:
      goal_stack = result.goals
      history.append((tactic_string, "success", result.goals))
  else:
      history.append((tactic_string, "error", result.error_message))
      # goal_stack unchanged — agent tries again from the same state,
      # now with the error message added to what it can see
  if budget_exhausted or agent_signals_stop:
      state = FAILED; stop
\end{verbatim}

Everything in this document is about making \texttt{lean.apply(...)} a real, callable thing, and about what \texttt{goal\_stack} and \texttt{result} actually look like on the wire --- because \emph{that} representation is what determines what the agent can see, and therefore what ``hallucination'' even means operationally (proposing a tactic that's invalid given what it could see, vs.\ one that's invalid period).

\section{Candidate interfaces, surveyed}

There is no single official ``Lean API for agents'' --- four real options exist, each with different tradeoffs:

\begin{center}
\begin{tabular}{@{}p{0.17\linewidth}p{0.24\linewidth}p{0.26\linewidth}p{0.27\linewidth}@{}}
\hline
\textbf{Interface} & \textbf{How it works} & \textbf{Strengths} & \textbf{Weaknesses} \\
\hline
\textbf{File + \texttt{lake env lean}} (from the Environment Setup chapter, \S8) & Write a whole \texttt{.lean} file, run \texttt{lake env lean file.lean}, parse stdout/stderr & Zero extra dependencies; exactly what CI and the setup doc already verify & No incremental state --- every attempt re-elaborates the whole file from scratch; painfully slow for a multi-step tactic loop; output is meant for humans, not machines \\
\textbf{leanprover-community/repl} & Persistent process, JSON in/out over stdin/stdout, \texttt{proofState} integers track state incrementally & Purpose-built for exactly this; incremental (no re-elaboration per tactic); structured JSON; supports backtracking via reusing an earlier \texttt{proofState}/\texttt{env} id & ``Experimental'' tactic mode per its own README; one Lean process = one project's environment, so multi-project setups need multiple processes (see \S5) \\
\textbf{Lean language server (LSP)}, via \texttt{lake serve} or wrapped as \texttt{lean-lsp-mcp} & The same server VS Code talks to; JSON-RPC 2.0; query goal state at a file position & Reuses infrastructure that's already battle-tested by every human Lean user; also exposes hover/definition/diagnostics beyond just tactic state & Designed around a \emph{live-edited file}, not a clean append-only tactic stream --- mapping ``propose one tactic'' onto ``edit this file at this position'' is more indirect than the REPL's native \texttt{\{"tactic": ..., "proofState": ...\}} shape \\
\textbf{Pantograph} & Purpose-built ML-facing API/REPL, supports tree search (MCTS) natively & Most feature-rich for search-heavy agent designs; used in published LLM+Lean pipelines (Draft-Sketch-Prove) & Heavier dependency, smaller community than the REPL above; overkill if the agent is doing simple sequential tactic proposal rather than tree search \\
\hline
\end{tabular}
\end{center}

\textbf{Recommendation for this project: the \texttt{leanprover-community/repl}, wrapped in Python (directly, or via \texttt{LeanInteract}).} It matches the project's stated loop shape almost exactly (see the tactic-mode example below), is incremental (fast --- no re-elaborating the whole file per tactic), and is simple enough to reason about and log precisely, which matters more here than raw proof-search power. Revisit Pantograph if/when the harness grows a tree-search or backtracking-heavy agent design rather than the linear one-tactic-at-a-time loop described in the project brief.

\section{The REPL protocol, concretely}

Run via \texttt{lake exe repl} (or, pointed at a specific project's dependencies: \texttt{lake env /path/to/repl/.lake/build/bin/repl}). Communicates via \textbf{one JSON object per command, blank-line separated}, on stdin, with one JSON object back on stdout.

\subsection{Starting a theorem (command mode to tactic mode)}

You open tactic mode by sending a command containing a \texttt{sorry}; the response gives you a \texttt{proofState} id to start from:

\begin{verbatim}
// → stdin
{"cmd": "theorem my_thm (x : Unit) : Nat := by sorry"}

// ← stdout
{"sorries":
 [{"proofState": 0,
   "pos": {"line": 1, "column": 29},
   "goal": "x : Unit\n⊢ Nat",
   "endPos": {"line": 1, "column": 34}}],
 "messages":
 [{"severity": "warning", "pos": {...}, "endPos": {...},
   "data": "declaration uses 'sorry'"}],
 "env": 0}
\end{verbatim}

\texttt{goal} here --- \texttt{"x : Unit\textbackslash{}n⊢ Nat"} --- is exactly the goal-state text this project's tutorial (the Lean 4 Tutorial chapter, \S1) teaches how to read. This is the string (or a structured equivalent) the agent should be shown at each step.

\subsection{One tactic step}

\begin{verbatim}
// → stdin
{"tactic": "apply Int.natAbs", "proofState": 0}

// ← stdout
{"proofState": 1, "goals": ["x : Unit\n⊢ Int"]}
\end{verbatim}

\begin{verbatim}
// → stdin
{"tactic": "exact -37", "proofState": 1}

// ← stdout
{"proofState": 2, "goals": []}
\end{verbatim}

\texttt{"goals": []} is the terminal success state --- the theorem is proved. This is the exact shape of \texttt{lean.apply()} from \S1: input is \texttt{\{tactic, proofState\}}, output is a new \texttt{proofState} plus the resulting \texttt{goals} list (empty on success).

\subsection{The error case}

An invalid tactic doesn't crash the process --- it returns a \texttt{messages} array with \texttt{"severity": "error"} and the proof state does \textbf{not} advance (you're still at the same \texttt{proofState}, free to try again):

\begin{verbatim}
// → stdin
{"tactic": "exact 37", "proofState": 1}   // wrong: goal wants Int, not the literal misapplied here

// ← stdout (schematic — exact message text depends on the specific mismatch)
{"proofState": 1,
 "messages": [{"severity": "error", "data": "type mismatch ..."}]}
\end{verbatim}

This is the critical mechanic for the hallucination analysis: \textbf{a rejected tactic does not lose progress} --- the agent is still at \texttt{proofState 1} and can retry. So ``how many iterations to converge'' (from the project's core question) is naturally just \emph{the count of tactic-submission round-trips}, separable into accepted-and-advanced vs.\ rejected-and-retried, per \texttt{proofState}.

\subsection{Backtracking and persistence}

Any earlier \texttt{env} or \texttt{proofState} integer can be reused as the starting point for a new command --- this \emph{is} backtracking, for free, without re-running anything. States can also be pickled to \texttt{.olean} files and restored later (\texttt{\{"pickleTo": "path", "proofState": n\}} / \texttt{\{"unpickleProofStateFrom": "path"\}}), which matters for \S6 (long-running logging) and for resuming an interrupted batch run.

\section{Data model: what gets exchanged per step}

Distilling the above into the fields the harness needs to track per iteration:

\begin{center}
\begin{tabular}{@{}p{0.24\linewidth}p{0.20\linewidth}p{0.50\linewidth}@{}}
\hline
\textbf{Field} & \textbf{Source} & \textbf{Notes} \\
\hline
\texttt{theorem\_id} & problem set (\texttt{problems/candidates.jsonl}, with its topic/difficulty split) & which problem this attempt belongs to \\
\texttt{attempt\_id} & harness & one per full proof attempt (a theorem may be attempted multiple times) \\
\texttt{step\_index} & harness & position in this attempt's tactic sequence \\
\texttt{proofState\_before} & REPL & the \texttt{proofState} id the tactic was proposed against \\
\texttt{tactic\_string} & agent output & exactly what the agent proposed, verbatim --- no normalization before logging \\
\texttt{proofState\_after} & REPL & new id if accepted; unchanged if rejected \\
\texttt{goals\_before} / \texttt{goals\_after} & REPL \texttt{goal}/\texttt{goals} field & the actual text shown to / returned to the agent \\
\texttt{outcome} & derived & \texttt{advanced} / \texttt{rejected} / \texttt{proved} (\texttt{goals == []}) / \texttt{no\_progress} (accepted syntactically but \texttt{goals\_after == goals\_before}) \\
\texttt{error\_message} & REPL \texttt{messages} & full text, only present on \texttt{outcome == rejected} \\
\texttt{wall\_time\_ms} & harness & time for this one REPL round-trip --- needed to separate ``slow tactic'' (e.g.\ a heavy \texttt{simp}/\texttt{aesop} call) from a hung process (\S7) \\
\texttt{llm\_input} & harness & the exact prompt sent to the model this step, as \texttt{\{system, user\}} --- the system prompt verbatim plus the user turn with the formatted goal and history \\
\texttt{llm\_raw\_response} & agent output & the completion exactly as returned, \emph{before} tactic extraction. Distinguishing this from \texttt{tactic\_string} is what separates ``the model proposed nonsense'' from ``our extractor mangled a good answer'' \\
\texttt{retry\_attempts} & harness & one entry per completion parsed this step, each carrying the raw text, the extracted tactic, and whether it was accepted or skipped as empty --- so an empty-completion retry is visible rather than inferred from a call count \\
\texttt{history\_passed} & harness & the history list as the model saw it at this step, logged before the step appends its own entry --- a true per-step snapshot, never back-filled \\
\hline
\end{tabular}
\end{center}

The first ten rows are the Lean-side record this document set out to specify.
The last four are the agent-side half that \S9 anticipated: the implemented
harness logs both, so a JSONL file is a complete transcript of the exchange
rather than a summary of it. Every field above ships in the per-step
\texttt{step} event, with \texttt{proofState}/\texttt{goals} naming following
the harness's own vocabulary (\texttt{goal\_before}, \texttt{goal\_after},
\texttt{lean\_messages}, \texttt{llm\_calls}, \texttt{lean\_calls}).

\section{Environment switching --- three distinct senses}

``Environment switching'' turns out to mean three different things at three different layers. Conflating them is an easy source of bugs once the harness is running many theorems across topics.

\subsection{Toolchain/project switching}

Each Lean \textbf{project} pins a \texttt{lean-toolchain} (exact compiler version) and a \texttt{lakefile} dependency set (which Mathlib commit, which other libraries). A single REPL process is started via \texttt{lake env <repl-binary>} \textbf{inside a specific project directory} and inherits \emph{that} project's toolchain and imports for its whole lifetime.

For this project's problem set (set theory / topology / algebra, all drawn from Mathlib per the Lean 4 Tutorial chapter, \S10--12), \textbf{one Mathlib-linked project should suffice for all three topics} --- Mathlib is a single library spanning all of them, so topic segmentation is a \emph{content/selection} concern, not an environment concern. Toolchain switching only becomes real if the problem set later pulls in an external benchmark repo with its own pinned toolchain (e.g.\ cloning \texttt{miniF2F} or \texttt{PutnamBench} directly rather than re-stating their theorems inside our own project). In that case, the harness needs one REPL process (and one working directory) per distinct toolchain, and a router that dispatches each theorem attempt to the correct process based on which benchmark it came from.

\subsection{Session/state switching}

Within a single REPL process, each theorem attempt starts from a fresh \texttt{env}/\texttt{proofState} (\S3). Retrying a theorem, or exploring two different tactic choices from the same point, is just reusing an earlier \texttt{proofState} id rather than starting over --- cheap, and exactly what makes backtracking-capable agent designs (or simply ``let me log both the accepted path and the road not taken'') tractable. The harness needs a clear policy here: does a new \emph{attempt} at the same theorem get a fresh \texttt{proofState 0}, or does it resume from wherever the last attempt got stuck? (Recommendation: fresh state per attempt, logged as a new \texttt{attempt\_id} --- resuming from a stuck point conflates ``the agent tried again'' with ``the agent kept going,'' which are different statistical events.)

\subsection{Interface switching (human vs.\ machine)}

This is the sense closest to the original phrasing of Goal 7 (the not-yet-written fluency guide) and to the Environment Setup chapter, \S8: VS Code + the Lean 4 extension is how \emph{you} build, debug, and sanity-check problems interactively (Infoview, hover, \texttt{exact?} suggestions rendered inline); the REPL (or LSP-via-MCP) is how the \emph{harness} talks to Lean headlessly, with no human in the loop. A theorem should generally be authored and manually verified once in VS Code (confirm it's actually provable, get a feel for its difficulty) before it's added to the problem set the harness runs against --- the REPL is not a good place to \emph{discover} whether a theorem statement itself is even correct, only to run an already-trusted theorem through repeated automated attempts.

\section{Failure modes at the plumbing level}

Distinct from the mathematical error taxonomy in the Lean 4 Tutorial chapter, \S13 (wrong term, wrong lemma name, missing instance, incomplete, no-progress, silent \texttt{sorry}), the harness also needs to handle failures in the \emph{plumbing itself} --- these are infrastructure bugs, not agent hallucinations, and conflating the two categories would corrupt the statistical analysis:

\begin{itemize}
  \item \textbf{Malformed JSON / protocol violation} --- the agent (or a bug in the harness) sends something the REPL can't parse. Should be caught and logged as a harness error, never attributed to the agent's mathematical reasoning.
  \item \textbf{Process crash} --- the REPL process dies mid-session (e.g.\ Mathlib bug, memory exhaustion). Needs a restart-and-resume policy; whatever attempt was in flight should be marked \texttt{incomplete\_infra\_failure}, not \texttt{failed}.
  \item \textbf{Non-terminating tactic} --- some tactics (\texttt{simp} with a bad lemma set, certain \texttt{aesop} configurations) can hang rather than error. The harness needs a \textbf{per-tactic timeout} with a hard kill, and that timeout event is itself a data point worth logging (which tactics the agent proposes that are prone to hanging is exactly the kind of pattern the hallucination analysis wants to catch).
  \item \textbf{Resource exhaustion across a long run} --- REPL processes accumulate \texttt{Environment}/\texttt{proofState} objects in memory; a long batch run needs periodic process restarts (using the pickling mechanism from \S3 to preserve any state worth keeping) rather than one REPL process for the entire experiment.
\end{itemize}

\subsection{What actually happened}

All four predicted failures materialised in practice. Three are fixed; one is not.

\begin{itemize}
  \item \textbf{Non-terminating tactic} --- fixed. \texttt{LEAN\_TACTIC\_TIMEOUT\_SECONDS = 120} bounds each tactic step with a hard kill. Elaboration gets its own, far longer budget (\texttt{LEAN\_ELABORATION\_TIMEOUT\_SECONDS = 900}), because a theorem submitted with its own \texttt{import Mathlib} header re-imports Mathlib inside the command --- measured at roughly 230\,s. Applying one flat timeout to both killed an entire 50-theorem run on its first theorem.
  \item \textbf{Process crash} --- fixed. \texttt{lean\_interact} kills the server when a command times out, so a single expired tactic would otherwise leave every later theorem facing a dead REPL that costs about ten minutes to rebuild. \texttt{\_ensure\_server\_alive()} restarts it, and both the elaboration and tactic call sites are wrapped so one bad theorem costs one theorem rather than the batch.
  \item \textbf{Malformed input / protocol violation} --- fixed, in the form it actually took. The REPL answers with either a structured response or a \texttt{LeanError} carrying only \texttt{message}; reading the latter as the former raised \texttt{AttributeError} and took down a whole run mid-batch. It is now classified as an ordinary error outcome.
  \item \textbf{Infrastructure versus reasoning} --- fixed, and load-bearing. The distinction this section insisted on is now carried by the result codes: \texttt{FAILED\_CONNECTIVITY\_OUTAGE} and \texttt{FAILED\_LLM\_ERROR} mark theorems the model never got to attempt, while \texttt{FAILED\_BUDGET\_EXHAUSTED} marks a genuine mathematical failure. Every pass rate in this study is reported twice because of it --- once over all theorems, once over only those that reached a working model.
  \item \textbf{Resource exhaustion across a long run} --- \emph{not} implemented. There are no periodic REPL restarts; a batch holds one process for its whole duration. The longest run to date kept a single REPL alive for roughly 42 hours without observable degradation, so this has not yet forced the issue --- but the prediction stands and the mitigation is still unwritten.
\end{itemize}

\section{Design decisions (resolved)}

These four were flagged as open when this document was written, on the grounds that they affect the statistics and so deserved deliberate choices rather than silent ones. All four are now settled in the implemented harness, and are recorded here so the analysis can state which variant was run:

\begin{enumerate}
  \item \textbf{What counts as ``one iteration''?} \emph{One LLM call is one step.} Each call produces exactly one \texttt{step} event in the JSONL log, whatever the proposed tactic contains --- a \texttt{<;>}-chained combinator counts once, because the agent took one turn. Every ``iterations to converge'' figure in the analysis is a count of these events.
  \item \textbf{Retry budget policy.} \emph{A fixed tactic budget, with full history visible.} \texttt{--max-steps} bounds the attempt (15 in the Easy-50 runs); there is no wall-clock budget, but a per-call timeout bounds each request. The agent sees its complete history every step --- every previous tactic, its outcome, and Lean's message --- and \texttt{history\_passed} records exactly what it was shown, so the claim is checkable rather than asserted.
  \item \textbf{What exactly gets shown to the agent on error.} \emph{The raw Lean error string, untruncated, identical across every run.} No cleaning, no summarisation, no per-run tuning. This was the choice the original text argued for, and it is worth keeping for the reason given there: a cleaner message could materially move the hallucination rate, so it has to be fixed and reported rather than adjusted.
  \item \textbf{Attempt independence.} \emph{Fresh proof state per attempt.} Each theorem starts from the \texttt{proofState} its own elaboration returned; a tactic is submitted against the current state, and nothing carries across attempts at the same theorem. Repeated attempts are therefore independent draws, matching \S5.2's recommendation.
\end{enumerate}

One consequence of (1) and (2) together is worth noting for the analysis: because an errored step does not advance the proof state, a 15-step budget is not 15 chances at 15 different goals --- it is 15 chances at however many distinct states the agent actually reached. In the Easy-50 runs the proved proofs closed within four steps and nothing closed later, so the budget was rarely the binding constraint.

\section{What this document fed --- and what changed}

Three deliverables were planned downstream of this chapter.
Two shipped; one proved unnecessary.

\textbf{Logging system (Goal 8).}
The agent-side schema described here shipped as four additional fields now recorded in every \texttt{step} event: \texttt{llm\_input} (the exact prompt sent), \texttt{llm\_raw\_response} (the raw completion before tactic extraction), \texttt{retry\_attempts} (a per-attempt trace of raw text, extracted tactic, and accepted/skipped status), and \texttt{history\_passed} (the tactic history the model received).
These are documented in \S4's event schema table.

\textbf{Problem set (Goal 9).}
\texttt{problems/candidates.jsonl} was populated with 837 theorems across three difficulty tiers (310 Easy, 378 Medium, 149 Hard) drawn from real Mathlib.
The planned hand-verification step --- authoring \texttt{.lean} files and checking elaboration in VS Code before any harness run --- was not completed.
Only 9 of the 837 theorems have hand-verified proofs; the remainder came from automated extraction.
The cost was visible: 4 of the first 50 Easy theorems failed elaboration outright (\texttt{FAILED\_ELABORATION}), which hand-verification would have caught.

\textbf{Environment-switching guide (Goal 7).}
Never written, and retrospectively unnecessary.
The harness runs a single Lean project, a single toolchain, and a single REPL process per batch.
The day-to-day switching problem the guide was meant to address did not materialise.
An unwritten deliverable that turned out to be unneeded is worth recording: the single-project decision argued for in \S5.1 resolved the problem the guide was meant to paper over.

\section{References}

\begin{itemize}
  \item \href{https://github.com/leanprover-community/repl}{leanprover-community/repl} --- protocol and JSON examples in \S3 quoted directly from this README
  \item \href{https://arxiv.org/abs/2410.16429}{Pantograph: A Machine-to-Machine Interaction Interface for Advanced Theorem Proving, High Level Reasoning, and Data Extraction in Lean 4} --- Aniva, Sun, Miranda, Barrett, Koyejo (arXiv:2410.16429)
  \item \href{https://github.com/oOo0oOo/lean-lsp-mcp}{lean-lsp-mcp}
  \item \href{https://github.com/augustepoiroux/leaninteract}{LeanInteract} --- Python wrapper; the implementation used throughout all runs
  \item The Lean 4 Tutorial chapter, \S13 (reading errors and diagnosing failed proofs), distinct from \S6's infrastructure failures here. The implemented failure taxonomy --- five modes: \textsc{No-Progress Loop}, \textsc{Error Loop}, \textsc{Hallucinated Lemma}, \textsc{Syntax Error Loop}, \textsc{Budget Exhaustion (Mixed)} --- lives in \texttt{classify\_failures.py}
  \item Agent-Lean systems this design was compared against: LeanDojo/ReProver, LeanAgent, Prover Agent, APOLLO
\end{itemize}
